% grammar.tex — Complete BNF Grammar for the Lain Programming Language
%
% This file is the single authoritative source for all grammar productions.
% Each rule carries a \label{gram:rule-name} for cross-referencing.
% Productions are grouped by clause; within each group, the order is
% top-down (start symbol first, auxiliary rules after).
%
% Usage in chapter files:
%   \gramref{expression}   → hyperlinked italic "expression"
%   To typeset a rule inline, use the grammar environment from lain-spec.sty.
%
% ★ Every production here was checked against the parser (src/frontends/lain/parser/) with a
% program for each form, 2026-10. The previous version had never been checked: it listed a
% block statement, `if` and block EXPRESSIONS, an `S { a: 1 }` constructor, a `->` token,
% `defer { }`, `proc`, `str`, and pointer types `*T[N]` and `*T[:S]`, none of which parse. It
% had no union type, no function-pointer type, no `for`, `break`, `continue`, `assert` or `case`
% statement, and no bitwise operators. Where the parser and the specification disagree, this
% file states the specification's rule, and the disagreement is a compiler defect: the
% precedence of `&` `^` `|` against the comparisons (chapter 8), a chained `as`, a `|` after a
% cast, a negative pattern after the first arm of a `case` statement, a negative bound in an
% alias refinement, and a leading zero in a decimal literal.
%
% ★ A listing does not wrap: keep every line inside a grammar environment within 88 columns,
% or the PDF clips it at the margin. Non-ASCII characters (§, —) and backquotes do not belong
% inside one either; a backquote prints as a curly quote.

% =========================================================================
% §5 — Lexical conventions
% =========================================================================

\section*{§5 Lexical grammar}

\subsection*{§5.3 Tokens}
\label{gram:token}
\begin{grammar}
token ::=
    keyword
  | identifier
  | integer-literal
  | float-literal
  | char-literal
  | string-literal
  | bool-literal
  | operator
  | punctuator

end ::= newline | ";"           % ends a statement or a declaration (chapter 5)

comment ::=
    "//" { any character but a line feed }
  | "/*" { any character | comment } "*/"           % block comments nest
\end{grammar}

\subsection*{§5.5 Keywords}
\label{gram:keyword}
\begin{grammar}
keyword ::=
    "and"       | "as"       | "assert"   | "assume"   | "break"
  | "c_include" | "case"     | "comptime" | "continue" | "decreasing"
  | "defer"     | "effects"  | "else"     | "extern"   | "false"
  | "for"       | "func"     | "if"       | "import"   | "in"
  | "mov"       | "nil"      | "or"       | "proc"     | "return"
  | "true"      | "try"      | "type"     | "unsafe"   | "use"
  | "var"       | "while"

% c_include, nil and proc are reserved and appear in no production: each was removed,
% and stays a keyword so that using it is diagnosed with the reason.
\end{grammar}

\subsection*{§5.6 Identifiers}
\label{gram:identifier}
\begin{grammar}
identifier           ::= identifier-start { identifier-continue }    % not a keyword
identifier-start     ::= letter | "_"
identifier-continue  ::= letter | decimal-digit | "_"
letter               ::= "A" .. "Z" | "a" .. "z"
qualified-identifier ::= identifier { "." identifier }
\end{grammar}

\subsection*{§5.7 Literals}
\label{gram:integer-literal}
\begin{grammar}
integer-literal  ::=
    decimal-literal
  | hex-literal
  | binary-literal
  | octal-literal

decimal-literal  ::= nonzero-digit { decimal-digit | "_" }
                   | "0"                          % no leading zero: 010 is ill-formed
hex-literal      ::= ( "0x" | "0X" ) hex-digit { hex-digit | "_" }
binary-literal   ::= ( "0b" | "0B" ) binary-digit { binary-digit | "_" }
octal-literal    ::= ( "0o" | "0O" ) octal-digit { octal-digit | "_" }

decimal-digit    ::= "0" .. "9"
nonzero-digit    ::= "1" .. "9"
hex-digit        ::= decimal-digit | "a" .. "f" | "A" .. "F"
binary-digit     ::= "0" | "1"
octal-digit      ::= "0" .. "7"
\end{grammar}

\label{gram:float-literal}
\begin{grammar}
float-literal    ::= digit-sequence "." digit-sequence [ float-exponent ]
                   | digit-sequence float-exponent
digit-sequence   ::= decimal-digit { decimal-digit | "_" }
float-exponent   ::= ( "e" | "E" ) [ "+" | "-" ] digit-sequence
\end{grammar}

\label{gram:string-literal}
\begin{grammar}
char-literal     ::= "'" ( char-element | escape-sequence ) "'"
char-element     ::= any character but "'", "\" or a line feed
string-literal   ::= '"' { string-element | escape-sequence } '"'
string-element   ::= any character but '"' or "\"
escape-sequence  ::= "\" ( "n" | "t" | "r" | "0" | "\" | "'" | '"' )
                   | "\" "x" hex-digit hex-digit
\end{grammar}

\label{gram:bool-literal}
\begin{grammar}
bool-literal     ::= "true" | "false"

literal          ::= integer-literal | float-literal | char-literal
                   | string-literal | bool-literal
\end{grammar}

\subsection*{§5.8 Operators and punctuators}
\label{gram:operator}\label{gram:punctuator}
\begin{grammar}
operator         ::=
    "+"   | "-"   | "*"   | "/"   | "%"
  | "+%"  | "-%"  | "*%"  | "/%"                   % wrapping
  | "+|"  | "-|"  | "*|"  | "/|"                   % saturating
  | "+?"  | "-?"  | "*?"                           % checked
  | "<<"  | "<<%" | ">>"
  | "&"   | "|"   | "^"   | "~"   | "!"
  | "=="  | "!="  | "<"   | ">"   | "<="  | ">="
  | "="   | "+="  | "-="  | "*="  | "/="  | "%="
  | "&="  | "|="  | "^="  | "<<=" | ">>="
  | "."   | ".."  | "..="
  | "and" | "or"  | "in"  | "as"

punctuator       ::=
    "("   | ")"   | "{"   | "}"   | "["   | "]"
  | ","   | ":"   | ";"   | "@"   | "?"   | "..."

% && and || are not operators; each is diagnosed with its spelling, and / or.
\end{grammar}

% =========================================================================
% §6 — Program structure and top-level declarations
% =========================================================================

\section*{§6 Program structure grammar}
\label{gram:program}\label{gram:import-declaration}
\begin{grammar}
program               ::= { top-level-declaration | end }

top-level-declaration ::=
    import-declaration
  | { attribute } type-declaration
  | { attribute } function-declaration
  | { attribute } extern-declaration
  | { attribute } constant-declaration
  | module-assertion

import-declaration    ::= "import" qualified-identifier
                          [ "." "{" identifier { "," identifier } "}" ]
                          [ "as" identifier ]
constant-declaration  ::= identifier [ type ] "=" expression     % compile-time constant
module-assertion      ::= "assert" expression             % checked at compile time

attribute             ::= "[" attribute-name [ "(" [ argument-list ] ")" ] "]" { end }
attribute-name        ::= "private" | "packed" | "ordered" | "fast_math"
                        | "cold" | "hot" | "allocator" | "noreturn"
                          % identifiers; any other name is E103
\end{grammar}

% =========================================================================
% §7 — Types
% =========================================================================

\section*{§7 Type grammar}
\label{gram:type}
\begin{grammar}
type             ::= core-type { "|" marker }       % a union: T | none, T | A | B
marker           ::= identifier
                     [ "(" marker-field { ( "," | newline ) marker-field } ")" ]
marker-field     ::= identifier type

core-type        ::=
    named-type { array-suffix }
  | "Vec" "(" integer-literal "," core-type ")"     % a SIMD vector of N lanes
  | "*" [ "var" ] core-type                         % raw pointer; *var T is writable
  | "*" "func" "(" [ core-type { "," core-type } ] ")" [ core-type ]
    [ effects-clause ]                              % function pointer
  | "mov" core-type                                 % owned
  | "var" core-type                                 % mutable borrow
  | "type"                                          % the type of a type parameter

% The pointee of * is not an array or slice type (E100): an array or a slice is
% already passed by reference.

named-type       ::= qualified-identifier
                     [ "(" type-argument { "," type-argument } ")" ]
type-argument    ::= core-type | integer-literal

% The primitive types are identifiers (chapter 7): int, i1 to i64, u1 to u64, isize,
% usize, float, f32, f64, bool, and void. A vector type may also be written as one
% identifier, element and lane count joined by x: u8x16 is Vec(16, u8).
\end{grammar}

\label{gram:array-type}
\begin{grammar}
array-suffix     ::=
    "[" integer-literal "]"                         % fixed-length array
  | "[" "]"                                         % slice
  | "[" [ refinement-op ] expression "]"            % slice, its length constrained
  | "[" ":" "0" "]"                                % a slice followed by a 0; any
                                                    % other sentinel is E100
\end{grammar}

\begin{grammar}
refinement       ::= range-refinement [ "and" comparison-list ]
                   | comparison-list
range-refinement ::= "in" bound-term ( ".." | "..=" ) bound-term
comparison-list  ::= refinement-op bound { "and" refinement-op bound }
refinement-op    ::= "<" | "<=" | ">" | ">=" | "==" | "!="
bound            ::= bound-term { ( "+" | "-" ) bound-term }
bound-term       ::= [ "-" ] integer-literal | identifier [ "." identifier ]

% Where each form is allowed (chapters 7, 10 and 12): a parameter takes all of them.
% A struct field's bound is one bound-term; a variant field's bound is a number. A
% return type's comparison-list has no negative number and no range. An alias's bound
% is a number or a name.
\end{grammar}

% =========================================================================
% §8 — Expressions
% =========================================================================

\section*{§8 Expression grammar}
\label{gram:expression}
\begin{grammar}
expression            ::= or-expression { "else" [ "return" ] or-expression }
                          % the error handler; it binds most loosely of all

or-expression         ::= and-expression { "or" and-expression }
and-expression        ::= equality-expression { "and" equality-expression }
equality-expression   ::= relational-expression
                          { ( "==" | "!=" ) relational-expression }
relational-expression ::= bitor-expression
                          { ( "<" | ">" | "<=" | ">=" ) bitor-expression
                          | "in" membership-set }
membership-set        ::= shift-expression ( ".." | "..=" ) shift-expression
                        | "[" expression { "," expression } "]"
bitor-expression      ::= bitxor-expression { "|" bitxor-expression }
bitxor-expression     ::= bitand-expression { "^" bitand-expression }
bitand-expression     ::= shift-expression { "&" shift-expression }
shift-expression      ::= additive-expression
                          { ( "<<" | "<<%" | ">>" ) additive-expression }
additive-expression   ::= multiplicative-expression
                          { additive-operator multiplicative-expression }
additive-operator     ::= "+" | "-" | "+%" | "-%" | "+|" | "-|" | "+?" | "-?"
multiplicative-expression ::= cast-expression
                          { multiplicative-operator cast-expression }
multiplicative-operator   ::= "*" | "/" | "%" | "*%" | "/%" | "*|" | "/|" | "*?"
cast-expression       ::= unary-expression { "as" [ "?" | "%" | "|" ] core-type }
unary-expression      ::= unary-operator unary-expression
                        | postfix-expression
unary-operator        ::= "-" | "!" | "~" | "&" | "*" | "mov" | "var" | "try"

postfix-expression    ::= ( identifier | "(" expression ")" ) { selector }
                        | primary-expression
selector              ::= "." identifier
                        | "(" [ argument-list ] ")"
                        | "[" index "]"
index                 ::= expression
                        | expression ".." [ expression ]
                        | ".." expression
argument-list         ::= expression { "," expression }

% A call, the construction of a struct (Point(1, 2)), of a variant (Shape.Circle(3))
% or of a marker with fields (Bad(3)), and an explicit generic argument
% (max(i32, a, b)) are one form: a name followed by "(" argument-list ")".

primary-expression    ::=
    literal
  | array-literal
  | array-comprehension
  | match-expression
  | builtin-expression

array-literal         ::= "[" [ expression { "," expression } [ "," ] ] "]"
array-comprehension   ::= "[" expression "for" identifier "in"
                          expression ( ".." | "..=" ) expression "]"

builtin-expression    ::=
    "@" ( "os" | "arch" )
  | "@" ( "likely" | "unlikely" ) "(" expression ")"
  | "@" ( "ctz" | "clz" | "popcount" | "movemask" ) "(" expression ")"
  | "@" ( "sizeof" | "alignof" ) "(" type ")"
  | "@" "splat" "(" type "," expression ")"
  | "@" "load" "(" type "," expression "," expression ")"
  | "@" "store" "(" expression "," expression "," expression ")"
  | "@" "shuffle" "(" expression "," expression ")"
  | "@" "assume_aligned" "(" expression "," integer-literal ")"
\end{grammar}

\label{gram:match-expression}
\begin{grammar}
match-expression      ::= "case" [ "&" ] expression "{" { expression-arm } "}"
expression-arm        ::= arm-head expression [ "," ]
arm-head              ::= ( pattern-list | "else" ) ":"
pattern-list          ::= pattern { "," pattern }
pattern               ::=
    pattern-value [ ( ".." | "..=" ) pattern-value ]   % a value, or a range of them
  | variant-name "(" [ identifier { "," identifier } ] ")"
                                                       % a variant, binding its fields
pattern-value         ::= [ "-" ] integer-literal | char-literal
                        | string-literal | bool-literal
                        | variant-name                 % a unit variant, or a constant
variant-name          ::= qualified-identifier
                        | identifier "(" type-argument { "," type-argument } ")"
                          "." identifier

% Arms are separated by newlines, or by ",". The body of an expression-arm is ONE
% expression; a "{" after the ":" is not a block. The statement form, whose arms hold
% statements, is match-statement (chapter 9).
\end{grammar}

% =========================================================================
% §9 — Statements
% =========================================================================

\section*{§9 Statement grammar}
\label{gram:statement}
\begin{grammar}
statement            ::=
    binding-declaration
  | variable-declaration
  | assignment-statement
  | expression-statement
  | if-statement
  | for-statement
  | while-statement
  | match-statement
  | "break"
  | "continue"
  | return-statement
  | defer-statement
  | unsafe-statement
  | assertion-statement
  | comptime-statement
  | use-statement

block                ::= "{" { end } [ statement { end { end } statement } ] { end } "}"

% There is no block STATEMENT: "{" never begins a statement. A block is the body of
% a function, an if or else, a loop, unsafe, or comptime if.

binding-declaration  ::= identifier [ type ] "=" expression             % immutable
variable-declaration ::= ( "var" | "mov" ) identifier [ type ]
                         [ "in" expression ] [ "=" expression ]
                         % var: mutable; mov: owned. Without "=", the type is required.

assignment-statement ::= assignable assignment-operator expression
assignable           ::= postfix-expression | "*" unary-expression
assignment-operator  ::= "=" | "+=" | "-=" | "*=" | "/=" | "%="
                       | "&=" | "|=" | "^=" | "<<=" | ">>="

expression-statement ::= expression

if-statement         ::= "if" expression block [ "else" ( if-statement | block ) ]
                         % else may begin the line after the "}"
for-statement        ::= "for" identifier [ "," identifier ] "in"
                         expression [ ( ".." | "..=" ) expression ] block
                         % for x in a, for i, x in a (index first), for i in lo..hi
while-statement      ::= "while" expression [ "decreasing" expression ] block

match-statement      ::= "case" [ "&" ] expression "{" { arm-head arm-body } "}"
arm-body             ::= { end } statement end { { end } statement end }
                         % statements, each ending at a newline, up to the next
                         % arm-head or the "}"

return-statement     ::= "return" [ expression ]
defer-statement      ::= "defer" statement               % one statement, not a block
unsafe-statement     ::= "unsafe" block
assertion-statement  ::= ( "assert" | "assume" ) expression
comptime-statement   ::= "comptime" "if" expression block [ "else" block ]
                       | "comptime" identifier [ type ] "=" expression
use-statement        ::= "use" qualified-identifier "as" identifier
                         % semantics implementation-defined (chapter 16); refused, E100
\end{grammar}

% =========================================================================
% §10 — Declarations
% =========================================================================

\section*{§10 Declaration grammar}
\label{gram:function-declaration}
\begin{grammar}
function-declaration ::= "func" identifier "(" [ parameter { "," parameter } ] ")"
                         [ return-type ] [ effects-clause ]
                         [ "decreasing" expression ] block

parameter            ::= [ "var" | "mov" ] identifier type [ refinement ]
                       | [ "var" | "mov" ] "{" identifier { "," identifier } "}" type
                         % the second form destructures the argument into its fields

return-type          ::= type [ "in" identifier ] [ comparison-list ]
                         % in p: the returned reference borrows parameter p

effects-clause       ::= "effects" [ effect-name { "," effect-name } ]
effect-name          ::= "io" | "diverge" | "raises" | "alloc"           % identifiers

extern-declaration   ::= "extern" "func" identifier "(" [ extern-parameters ] ")"
                         [ type [ "in" identifier ] ] [ effects-clause ]
                       | "extern" "type" identifier               % an opaque C type
extern-parameters    ::= identifier type { "," identifier type } [ "," "..." ]
                       | "..."
\end{grammar}

\label{gram:struct-declaration}\label{gram:enum-declaration}
\begin{grammar}
type-declaration     ::= alias-declaration | struct-declaration | enum-declaration

alias-declaration    ::= "type" identifier "=" alias-target
alias-target         ::=
    identifier comparison-list           % a refined alias: type Small = u8 < 200
  | named-type { "[" expression "]" }    % a constant length: type Grid = u8[3][3]
  | "type" "{" ( field-list | variant-list ) "}"   % an anonymous struct or enum

struct-declaration   ::= "type" identifier [ type-parameters ] "{" [ field-list ] "}"
field-list           ::= { end } field { ( "," | end ) { end } field }
                         [ "," ] { end }
field                ::= [ "mov" ] identifier type [ refinement ]

enum-declaration     ::= "type" identifier [ type-parameters ] [ backing-width ]
                         "{" [ variant-list ] "}"
variant-list         ::= { end } variant { ( "," | end ) { end } variant }
                         [ "," ] { end }
variant              ::= identifier
                         [ "{" variant-field { ( "," | end ) variant-field } "}" ]
variant-field        ::= identifier type [ comparison-list ]
                         % a linear field is h mov H: here mov belongs to the type
backing-width        ::= "u8" | "u16" | "u32" | "u64"          % a plain enum only

type-parameters      ::= "(" identifier type { "," identifier type } ")"
                         % T type is a type parameter; N usize a compile-time value
\end{grammar}
